% Requires: amsmath, booktabs, algorithm, algpseudocode, enumitem
\section{System Architecture}
\label{sec:architecture}

HybridCloudSim is a discrete-event simulator built on SimPy. It models a hybrid
cloud at the \emph{orchestration} layer: job arrival, device selection,
capacity blocking, iterative quantum--classical execution, and the resulting
energy and cost. Circuit-level quantum behavior is deliberately abstracted;
device service times are obtained from analytical performance models
parameterized by published hardware specifications, so that a full run over
thousands of jobs completes in seconds on a single core.

\subsection{Overview}
\label{sec:arch-overview}

The simulator is organized around a single composition root,
\texttt{HybridCloudSimEnv}, which subclasses the SimPy environment and
instantiates every other component. Table~\ref{tab:components} summarizes the
roles. Control flows top-down --- the job feed creates one broker process per
job, and each broker drives devices --- while observability flows bottom-up:
devices emit events on a publish/subscribe bus and write per-job records that
the monitor and the records manager aggregate.

\begin{table}[t]
\centering
\caption{Principal components of HybridCloudSim.}
\label{tab:components}
\small
\begin{tabular}{@{}ll@{}}
\toprule
\textbf{Component} & \textbf{Responsibility} \\
\midrule
\texttt{HybridCloudSimEnv} & Composition root; owns the cost model \\
\texttt{JobGenerator}      & Trace replay or stochastic arrivals \\
\texttt{HybridBroker}      & Per-job scheduling and phase control \\
\texttt{QuantumDevice}     & Qubit capacity, topology, QPU service time \\
\texttt{CPU} / \texttt{AMDRyzen} & Core and memory-bandwidth capacity \\
\texttt{EventBus}          & Decoupled device event notification \\
\texttt{CloudMonitor}      & Utilization integration, fleet power \\
\texttt{JobRecordsManager} & Event log, derived metrics, energy and cost \\
\bottomrule
\end{tabular}
\end{table}

A run is fully specified by three inputs: a device inventory (a list of QPU and
CPU objects), a workload (a trace file or an arrival process), and a cost
configuration. Devices are constructed detached from the environment and are
bound to it in a separate initialization pass, at which point their SimPy
containers and resources are created. This two-phase construction keeps device
definitions declarative and makes parameter sweeps explicit: because devices
carry mutable state --- qubit containers, the topology graph, and the qubit
occupancy map --- a fresh inventory is instantiated for each configuration in a
sweep.

\subsection{Workload Model}
\label{sec:arch-workload}

A job is a tuple
$J=\langle n, d, s, p, a, r, c, b\rangle$
giving qubit count $n$, circuit depth $d$, shot count $s$, priority $p$,
arrival time $a$, required iterations $r$, CPU units $c$, and memory bandwidth
$b$. Jobs enter the system in one of two modes. In \emph{dispatcher} mode the
generator replays a CSV or JSON trace, releasing each job at its recorded
arrival time; this is the mode used for all experiments in this paper, since it
makes workloads reproducible and shareable. In \emph{generator} mode jobs are
synthesized online from a configurable inter-arrival distribution
(exponential by default) with randomized circuit parameters. In both modes the
generator spawns one broker process per job, so concurrency in the model is
bounded by device capacity rather than by any fixed pool of workers.

\subsection{Broker and Execution Lifecycle}
\label{sec:arch-broker}

The broker is the scheduling core. Each job executes $r$ variational
iterations, and every iteration consists of a QPU phase followed by a CPU
phase, reflecting the structure of VQE- and QAOA-style workloads in which
measurement outcomes are returned to a classical optimizer that produces the
next parameter set. Algorithm~\ref{alg:broker} states the loop.

\begin{algorithm}[t]
\caption{Broker process for job $J$}
\label{alg:broker}
\begin{algorithmic}[1]
\State $i \gets 0$
\While{$i < r$}
  \State $q \gets \Call{Select}{\text{QPU}, n}$
  \While{$q = \varnothing$} \Comment{no device has free capacity}
    \State \textbf{wait} $\delta$; \quad $q \gets \Call{Select}{\text{QPU}, n}$
  \EndWhile
  \State \Call{Execute}{$q, J$}; \quad \Call{Record}{$J$, QPU}
  \State $u \gets \Call{Select}{\text{CPU}, (c,b)}$
  \While{$u = \varnothing$}
    \State \textbf{wait} $\delta$; \quad $u \gets \Call{Select}{\text{CPU}, (c,b)}$
  \EndWhile
  \State \Call{Execute}{$u, J$}; \quad \Call{Record}{$J$, CPU}
  \State $i \gets i + 1$
\EndWhile
\State \Call{FinalizeEnergyAndCost}{$J$}
\end{algorithmic}
\end{algorithm}

\textsc{Select} filters the inventory to devices of the requested type that are
not under maintenance and whose free capacity meets the request --- free qubits
for a QPU, free cores \emph{and} free memory bandwidth for a CPU --- and
returns the most lightly loaded candidate, breaking ties in favor of the
device with the largest residual capacity. When no candidate exists the broker
retries after a fixed polling interval $\delta$ rather than enqueueing, which
keeps the admission decision re-evaluated against current fleet state. The
resulting blocking behavior, rather than an explicit queue discipline, is what
produces the contention observed in Section~\ref{sec:evaluation}. The policy is
localized in \textsc{Select}, so alternative disciplines (priority, fair share,
energy-aware placement) can be substituted without touching the phase loop.

\subsection{Two-Level QPU Allocation}
\label{sec:arch-qpu}

Quantum devices are the only resource in the model with a spatial constraint.
Each device loads a topology graph $G=(V,E)$ from a hardware description file,
where $|V|$ is the physical qubit count, and maintains an occupancy map marking
each qubit free or allocated. Admission proceeds at two levels. First, an
aggregate counter must hold at least $n$ free qubits. Second, the device must
contain a \emph{connected} induced subgraph of $n$ free qubits, located by a
breadth-first search over free vertices. On allocation the selected qubits are
marked busy and their edges to the remainder of the graph are cut; on
completion the qubits are released and the cut edges are restored between pairs
that are both free again.

This second level is what distinguishes the model from capacity-only
abstractions. A job can satisfy the aggregate test and still fail to be
admitted because the free qubits are fragmented across disconnected regions of
the lattice, which is the expected behavior on sparsely connected
superconducting topologies under multi-tenant load.

QPU service time follows a throughput model based on the vendor-reported
Circuit Layer Operations Per Second (CLOPS) metric:
\begin{equation}
\label{eq:qpu-time}
t_{\mathrm{QPU}} \;=\; \alpha\,\frac{M K S \log_2 d}{\mathrm{CLOPS}},
\end{equation}
where $S$ is the shot count, $d$ the circuit depth, $M$ and $K$ the number of
circuit templates and parameter updates per submission, and $\alpha$ a unit
conversion constant. Per-device calibration tables (readout, single-qubit, and
two-qubit gate errors) are additionally loaded to expose an estimated circuit
fidelity alongside timing, though fidelity does not feed back into scheduling
in the present study. The library ships topology and calibration data for
sixteen IBM superconducting processors as well as D-Wave, Rigetti, and Google
devices; adding a device requires only a topology file and its specification
constants.

\subsection{Classical Device Model}
\label{sec:arch-cpu}

Classical nodes expose two independent capacities, compute units and memory
bandwidth, both of which must be satisfied for a job to be admitted; modeling
bandwidth separately is necessary because statevector post-processing is
frequently bandwidth- rather than compute-bound. The default node derives its
service time from the algorithmic cost of the classical half of the
variational loop:
\begin{equation}
\label{eq:cpu-time}
t_{\mathrm{CPU}} \;=\; \frac{\kappa\, 2^{\,n} d\, \pi}{\eta\, c^{\,0.85}},
\qquad
\eta = \tfrac{1}{2}\bigl(f_{\mathrm{base}} + f_{\mathrm{boost}}\bigr)\cdot\mathrm{IPC},
\end{equation}
where $\pi$ is the variational parameter count, $\kappa$ a calibration
constant, $\eta$ the effective per-core throughput derived from clock and IPC,
and the exponent $0.85$ captures sublinear scaling with allocated units. The
$2^{\,n}$ term makes the classical side of the loop grow exponentially with
qubit width, which is the mechanism by which classical post-processing
overtakes quantum execution as circuits widen. A simpler node with randomized
service time is also provided for baseline comparisons.

\subsection{Energy and Cost Model}
\label{sec:arch-energy}

Power parameters live in a single configuration dictionary passed to the
environment, with per-device overrides; all figures reported in
Section~\ref{sec:evaluation} are obtained by varying this dictionary alone.
Superconducting QPUs are modeled as drawing a constant power $P^{q}$ dominated
by the cryogenic plant whenever the device hosts any job, since dilution
refrigeration is essentially load-independent. Classical nodes follow the
CloudSim affine model, interpolating between idle and peak draw by utilization
$u_j$. Instantaneous fleet power is therefore
\begin{equation}
\label{eq:power}
P(t) \;=\; \sum_{q \in \mathcal{Q}} P^{q}\,\mathbb{1}\!\left[\text{$q$ busy at $t$}\right]
\;+\; \sum_{j \in \mathcal{C}} \Bigl[P^{j}_{\mathrm{idle}} + \bigl(P^{j}_{\mathrm{peak}} - P^{j}_{\mathrm{idle}}\bigr) u_j(t)\Bigr].
\end{equation}
Per-job energy is attributed by phase. For job $J$ with phase durations
$t^{q}_{i}$ and $t^{c}_{i}$ on the devices selected in iteration $i$,
\begin{equation}
\label{eq:energy}
E(J) \;=\; \sum_{i=1}^{r}\Bigl(P^{q}_{i} t^{q}_{i} + P^{c}_{i} t^{c}_{i}\Bigr),
\qquad
\mathrm{Cost}(J) \;=\; \rho\, E(J),
\end{equation}
with $\rho$ the electricity price per kWh. Equations~\eqref{eq:power}
and~\eqref{eq:energy} are complementary rather than redundant: the former gives
the operator's instantaneous load, which determines provisioning and peak
demand charges, while the latter gives the per-tenant attribution needed to set
prices. Both are driven by the same configuration, so a change to a device's
power profile propagates to both views.

\subsection{Telemetry}
\label{sec:arch-telemetry}

Devices publish \texttt{device\_start} and \texttt{device\_finish} events on a
lightweight event bus. The cloud monitor subscribes to both and, on each event,
integrates the resource-time product accumulated since the previous event
before recording a snapshot containing time-averaged utilization for qubits,
cores, and bandwidth together with instantaneous power from
Equation~\eqref{eq:power}. Because utilization is piecewise constant between
allocation events, event-boundary sampling is exact rather than approximate,
and the sample count scales with system activity instead of with wall-clock
duration.

In parallel, the records manager maintains an append-only event log per job.
Each phase contributes arrival, start, and finish stamps, from which waiting
time, service time, and turnaround are derived per phase and per iteration;
makespan is computed on the final iteration, at which point
Equation~\eqref{eq:energy} is evaluated and the energy and cost fields are
written back. Every field is retained as a per-iteration sequence, so the
traces support both distributional analysis across jobs and within-job analysis
across iterations. The two telemetry paths --- integrated fleet-level and
discrete per-job --- are reconciled in Section~\ref{sec:evaluation}, where
aggregate energy from the job records is compared against the time integral of
the monitored power series.
